\subsection{可分闭域上的代数}

引理 1. —— 设 D 是一个体，Z 是它的中心。用 p 表示 Z 的特征指数。假设对 D 的每个元素 a，存在整数 \(m \geqslant 0\) 使得 \(a^{p^{m}}\) 属于 Z。则体 D 是交换的。

若 p = 1，则 D = Z。因此我们假设 p \textgreater{} 1。

我们用反证法来证明，假设 D 不是交换的。设 \(a\) 是 \(\mathrm{D} - \mathrm{Z}\) 的一个元素，\(q\) 是 \(p\) 的一个幂，使得 \(a^q\) 属于 Z。用 I 表示 D 上的恒等映射，把内自同构 \(x \mapsto axa^{-1}\)（从 D 到 D、与 \(a\) 相伴）记作 \(\sigma\)；我们有 \(\sigma^q = \mathrm{I}\)，因为 \(a^q\) 属于 Z。我们有 \(\sigma - \mathrm{I} \neq 0\)，因为 \(a\) 不属于 Z；我们还有 \((\sigma - \mathrm{I})^q = \sigma^q - \mathrm{I} = 0\)，因为 Z 的特征为 \(p\)。设 \(f\) 是使 \((\sigma - \mathrm{I})^f \neq 0\) 的最大自然数；我们有 \(f \geqslant 1\)。选取 D 的一个元素 \(c\)，使得 \((\sigma - \mathrm{I})^f(c) \neq 0\)，并令

\[
x = (\sigma - \mathrm{I}) ^ {f - 1} (c), \qquad y = (\sigma - \mathrm{I}) (x) = (\sigma - \mathrm{I}) ^ {f} (c).
\]

按构造，我们有 \(y \neq 0\) 且 \(\sigma(y) = y\)；若令 \(z = y^{-1}x\)，则由此得出

\[
\sigma (z) = \sigma (y) ^ {- 1} \sigma (x) = y ^ {- 1} (y + x) = 1 + z
\]

因此 \(\sigma(z^{p^j}) = 1 + z^{p^j}\) 对每个自然数 \(j\) 成立。选取整数 \(m \geqslant 0\) 使得 \(z^{p^m}\) 属于 D 的中心 Z；我们有

\[
z ^ {p ^ {m}} = a z ^ {p ^ {m}} a ^ {- 1} = \sigma (z ^ {p ^ {m}}) = 1 + z ^ {p ^ {m}}.
\]

这一矛盾便导出引理 1。

命题 3. —— 设 K 是一个可分闭域（V, §7, No.~8, 第 45 页，定义 4），D 是 K 上的一个为体的有限次数代数。则 D 是交换的。

用 p 表示 K 的特征指数。设 a 是 D 的一个元素。环 \(K[a]\) 是 K 的一个代数扩张（V, §3, No.~1, 第 17 页，推论 1）。由于域 K 是可分闭的，由 V, §7, No.~7, 第 44 页的命题 13 得出，代数 \(K[a]\) 是 K 的一个 p-根式扩张。于是存在整数 \(m \geqslant 0\) 使得 \(a^{p^{m}}\) 属于 K。由引理 1，体 D 因而是交换的。

推论. —— 设 K 是一个可分闭域，A 是 K 上的一个半单的有限次数代数。则存在整数 \(r \geqslant 0\)、严格正整数 \(n_{1}, \ldots, n_{r}\) 以及 K 的有限次数扩张 \(K_{1}, \ldots, K_{r}\)，使得 A 同构于代数 \(\prod_{i=1}^{r} \mathbf{M}_{n_{i}}(\mathbf{K}_{i})\)。

由半单代数的结构定理（VIII, 第 135 页，定理 1），A 同构于一个代数 \(\prod_{i=1}^{r} \mathbf{M}_{n_i}(\mathrm{D}_i)\)，其中 \(r\) 是整数 \(\geqslant 0\)，\(n_1, \ldots, n_r\) 是严格正整数，而 \(\mathrm{D}_1, \ldots, \mathrm{D}_r\) 是 K 上的为体的有限次数代数。由于域 K 是可分闭的，由命题 3，每个体 \(\mathrm{D}_i\) 都是交换的；推论得证。

